\documentclass{article}
\usepackage{../../math_template}

\author{Jonathan Kasongo}
\title{Noether Sheet 1}

\begin{document}
\maketitle

\noindent
These were my solutions to the Noether Mentoring Scheme Sheet 1.
Problem statements are abridged, and
\centerline{problem 1 is omitted because it isn't
a math problem.}

\begin{problem}{Problem 2}
17:36 29/05/84 is a \textit{pandigital} time and date as it uses each of
the digits 0-9 exactly once. \vspace{0.2cm}

What is the first pandigital time and date of each century?
\end{problem}

\begin{solution}{Solution to problem 2}
The day cannot be 30 or 31. Suppose the contrary, the tens digit of the
hour is 0,1 or 2 and the tens digit of the month is 0 or 1. Thus the
hour must use the number 2, so that the month can use 0 if the day is 31 and
1 if the day is 30. So now the numbers 0,1,2 and 3 are all taken. In the
hour the only numbers that can follow that 2 are 0,1,2 and 3, but since all
of them are taken we will repeat a digit, contradiction.
\vspace{0.2cm}

Now we notice that these positions can only contain numbers from the set
$\{0, 1, 2\}$ (the tens digit of the month can't be 2 though).

\large
$$
\text{
\underline{*} \underline{\phantom{a}} : \underline{\phantom{a}} \underline{\phantom{a}} \ \ \
\underline{*} \underline{\phantom{a}} / \underline{*} \underline{\phantom{a}} / \underline{\phantom{a}} \underline{\phantom{a}}
}
$$

\normalsize
Also if we make the tens digit of the hour a 2 then it must be followed by
a 3 since 0,1 and 2 are taken, but if we make the tens digit of the hour a
1 then we can freely put whatever digit after it, allowing us to use
the digit 3 in the tens digit of the year.
This is optimal 3 is the smallest digit available to use.
Finally if we make the tens digit of the hour a 0 then the month will be
$\overline{1x} > \overline{0x}$ so this is unoptimal. Putting this
information together we have

\large
$$
\text{
\underline{1} \underline{\phantom{a}} : \underline{\phantom{a}} \underline{\phantom{a}} \ \ \
\underline{2} \underline{\phantom{a}} / \underline{0} \underline{\phantom{a}} / \underline{\phantom{a}} \underline{\phantom{a}}
}
$$
\normalsize

Now we proceed greedily, it is optimal to make the year 34. We can't make
the month 05 because otherwise the tens digit of the minutes would cause
a repitition (the numbers 0,1,2,3,4 and 5 would be taken), so it is best to
make the month 06. As reasoned the tens digit of the minutes must be 5 since
there are no other options. The time and date now look like so

\large
$$
\text{
\underline{1} \underline{\phantom{a}} : \underline{5} \underline{\phantom{a}} \ \ \
\underline{2} \underline{\phantom{a}} / \underline{0} \underline{6} / \underline{3} \underline{4}
}
$$
\normalsize

Now there are no more restrictions on the digits we can just greedily fill
in the blanks and the final answer comes out as

\large
$$
\text{
\underline{1} \underline{8} : \underline{5} \underline{9} \ \ \
\underline{2} \underline{7} / \underline{0} \underline{6} / \underline{3} \underline{4}
}
$$
\normalsize

$\blacksquare$\\

\textit{Comments: This problem looks deceptively easy. I fakesolved many
times, I wonder if there is a more elegant approach here. }
\end{solution}
\vspace{0.5cm}
\newpage

\begin{problem}{Problem 3}
Derive the quadratic formula for $ax^2 + bx + c$ given $b^2 \geq 4ac$.
\end{problem}

\begin{solution}{Solution to problem 3}
Complete the square,

\[
\begin{aligned}
a\left(x+\frac{b}{2a}\right)^2 &= \frac{b^2 - 4ac}{4a} \\
\sqrt{a}\left(x+\frac{b}{2a}\right) &= \pm\frac{1}{2\sqrt{a}}\sqrt{b^2-4ac}
\\
x &= -\frac{b}{2a} \pm \frac{1}{2a}\sqrt{b^2 - 4ac}\\
\end{aligned}
\]

as required. $\blacksquare$\\

\textit{Comments: This is a pretty standard problem and I've seen the
derivation before so it was solved quickly.}
\end{solution}
\vspace{0.5cm}
\newpage

\begin{problem}{Problem 4}
A triangle has sides of length 13 cm, 14 cm and 15 cm.

\begin{enumerate}[label=(\alph*)]
\item Show that this triangle can be split into two right angled triangles with integer side lengths (when
measured in centimetres).
\item A rectangle is drawn inside the triangle, with its base on one of the sides of the triangle.
What is the largest possible area of this rectangle?
\end{enumerate}

\end{problem}

\begin{solution}{Solution to problem 4}
\textbf{(a):} Let $\triangle ABC$ be a triangle such that $AB = 15$ cm,
$CA = 14$ cm and $BC = 13$ cm. Let $F$ be the foot of the altitude that
passes through $B$. Let $FC = a$ cm and $FA = b$ cm, and notice that
$\triangle BFC$ and $\triangle BFA$ are both right angled triangles.
Let $BF = x$ cm.


\begin{center}
\begin{asy}
size(120);

pair A = (0,0);
pair C = (14,0);
pair F = (9,0);
pair B = (9,12);

pen pastel_blue = rgb(0.68, 0.85, 0.90);
pen pastel_pink = rgb(0.99, 0.71, 0.75);
pen line_pen = black + linewidth(1);

fill(A--F--B--cycle, pastel_blue + opacity(0.2));
fill(F--C--B--cycle, pastel_blue + opacity(0.2));

draw(A--B, blue);
draw(B--C, blue);
draw(C--A, blue);
draw(B--F, blue);

label("$A$", A, SW);
label("$C$", C, SE);
label("$B$", B, N);
label("$F$", F, S);

label("$15$", midpoint(A--B), NW);
label("$13$", midpoint(B--C), NE);
label("$b$", midpoint(A--F), S);
label("$a$", midpoint(F--C), S);
label("$x$", midpoint(B--F), E);
\end{asy}
\end{center}


We have a system of equations.

\[
\begin{aligned}
a^2 + x^2 &= 13^2\\
b^2 + x^2 &= 15^2
\end{aligned}
\]

and so $b^2 - a^2 = (b-a)(14) = 56$ since $b+a = 14$. So we deduce that
$b-a = 4$, from which we can solve to see that $b = 9, a = 5$ and also
$x = \sqrt{13^2 - 5^2} = 12$. Thus the triangles $\triangle BFC$ and
$\triangle BFA$ are both right-angled and also have integer side lengths.
$\blacksquare$ \\

\textbf{(b):} We begin with a claim. \vspace{0.2cm}


\textit{Claim 1: For any non-degenerate triangle $\triangle ABC$, let
$R$ be a rectangle inscribed within $\triangle ABC$ such that one of it's
bases lies on one of the sides of $\trinangle ABC$. Then
$\text{Area}(R) \leq \frac12 \text{Area}(\triangle ABC)$.
} \vspace{0.2cm}

\textit{Proof:} Let the base of $R$ lie on some side $P$ of
$\triangle ABC$ that has length $p$. Let the height of
$\triangle ABC$ be $h$ (modelling $P$ to be the base of the triangle).
Let $R$ have vertices $r_1, r_2, r_3$ and $r_4$, (starting from the top
left and going clockwise). Let the height of $R$ be $(1-\lambda)h$ and let
the length of $R$ be $\mu p$ for some $0< \mu, \lambda <1$.
The following diagram can help us visualise,

\begin{center}
\begin{asy}
/* Geogebra to Asymptote conversion, documentation at artofproblemsolving.com/Wiki go to User:Azjps/geogebra */
import graph; size(120);
real labelscalefactor = 1.0; /* changes label-to-point distance */
real xmin = -9.809343418249249, xmax = 4.164122123731424, ymin = -2.0835419068351797, ymax = 7.156476636053874;  /* image dimensions */
pen ttttff = rgb(0.2,0.2,1.); pen zzttqq = rgb(0.6,0.2,0.); pen ccqqqq = rgb(0.8,0.,0.);
pair A = (-5.,0.), B = (3.,0.), C = (-3.,6.), r_1 = (-4.,3.), r_2 = (0.,3.), r_3 = (0.,0.), r_4 = (-4.,0.);
pen pastel_blue = rgb(0.68, 0.85, 0.90);


filldraw(A--B--C--cycle, pastel_blue + opacity(0.2) ,ttttff);
filldraw(r_1--r_2--r_3--r_4--cycle, zzttqq + opacity(0.2), zzttqq + opacity(0));
/* draw figures */
draw(A--B, ttttff);
draw(B--C, ttttff);
draw(C--A,  ttttff);
draw(r_1--r_2, dashed + ccqqqq);
draw(r_2--r_3, dashed + ccqqqq);
draw(r_3--r_4, zzttqq);
draw(r_4--r_1, dashed + ccqqqq);
draw(C--(-3.,3.), arrow=Arrows);
draw((-2.,6.)--(-2.,0.), arrow=Arrows);
/* dots and labels */
label("$A$", (-4.931521129424998,0.10271248054481974), NE * labelscalefactor);
label("$B$", (3.040057367955859,0.10271248054481974), NE * labelscalefactor);
label("$C$", (-2.9618296766439585,6.1045995251446294), NE * labelscalefactor);
label("$p$", (-1.0024507445581072,-0.85510053683489716), NE * labelscalefactor,ttttff);
label("$r_1$", (-3.900269059906129,3.093343482149536), NE * labelscalefactor);
label("$r_2$", (0.039113845655950406,3.083030961454347), NE * labelscalefactor);
label("$r_3$", (0.039113845655950406,0.10271248054481974), NE * labelscalefactor);
label("$r_4$", (-3.9621441840772613,0.10271248054481974), NE * labelscalefactor);
label("$\mu p$", (-2.00276525199141,0.17490012541114047), NE * labelscalefactor,zzttqq);
label("$\lambda h$", (-3.168080090547732,4.506158817390385), NE * labelscalefactor);
label("$h$", (-2.1677655831144294,3.3005307958928378), NE * labelscalefactor);
clip((xmin,ymin)--(xmin,ymax)--(xmax,ymax)--(xmax,ymin)--cycle);
/* end of picture */
\end{asy}
\end{center}

Using corresponding angles and the fact that $\angle ACB$ is common to
both triangles we deduce that $\triangle ABC$ and
$\triangle C r_1 r_2$ are similar, thus it is the case that
$\lambda h / h = \mu p / p$. Thus it is the case that
$\mu = \lambda$. The area of $R$ is $\mu p \cdot (1-\lambda)h =
\mu (1-\mu) ph$. It is easy to check that the maxima of $f(x)=x(1-x)$
is $\frac14$ and so we find that

$$
\text{Area}(R) \leq \frac14 ph =
\frac12 \text{Area}(\triangle ABC).\ \square
$$



Now due to part (a) we know that $\text{Area}(\triangle ABC) =
\frac12 (12)(14)=84$, and so by claim 1 the largest possible area of the
rectangle is 42. $\blacksquare$
\\

\textit{Comments: I think this problem was my favourite just after problem
8. The motivation for conjecturing
$\text{Area}(R) \leq \frac12 \text{Area}(\triangle ABC)$, came from
considering a rectangle with base on the side $AC$ using the triangle
in part (a). One can form a quadratic for the area of the rectangle in
terms of the height of the
rectangle which attains a maxima of 42. Now that answer is
fairly suspicious since $42 = \frac12 84 =
\frac12 \text{Area}(\triangle ABC)$, so from here it is natural to
guess that claim 1 might be true.
}
\end{solution}
\vspace{0.5cm}
\newpage

\begin{problem}{Problem 5}
\begin{enumerate}[label=(\alph*)]
\item Aadya cycles down a hill at $x$ m/s and up the same hill at $y$ m/s. Barry then estimates her average
speed by calculating $\frac{x+y}{2}$ m/s.
Show that Barry’s estimation is always greater than or equal to Aadya’s true average speed.
\item One day, Aadya cycled down the hill at $x$ m/s and up the hill at 3 m/s. Both $x$ and the difference
between Barry’s estimation of Aadya’s average speed and her true average speed are integers. What are the possible values of $x$?
\end{enumerate}
\end{problem}

\begin{solution}{Solution to problem 5}
\textbf{(a):} Suppose that the hill is $h$ metres long. For some real
positive times $t_1,t_2$ we have
$
x = h/t_1,\ y = h/t_2.
$

The average speed of Aadya is
$$
\overline{s}:=\frac{2h}{t_1 + t_2} = \frac{2xy}{x+y}
$$

Now consider

\[
\begin{aligned}
\frac{2xy}{x+y} &\leq \frac{x+y}{2} \\
4xy &\leq (x+y)^2 \\
2xy &\leq x^2 + y^2 \\
\end{aligned}
\]


but this is true by AM-GM. $\blacksquare$
\\

\textbf{(b):} Similarly we write $x=h/t_1,\ 3=h/t_2$. The average speed of
Aadya is
$$
\overline{s} := \frac{2h}{t_1+t_2} = \frac{6x}{x+3}
$$

Now we are told that

$$
\frac{x+3}{2} - \frac{6x}{x+3} = \frac{(x-3)^2}{2(x+3)} \in
\mathbb{Z}_{\geq 0}
$$

since Barry's estimation is always greater than or equal to Aadya'a true
speed. We are asking to find all $x \in
\mathbb{Z}_{\geq 0}$ such that $2(x+3) \mid (x-3)^2$. If $x$ were even
then the numerator isn't divisible by 2 whilst the denominator is divisible
by 2, a contradiction. So $x$ must be odd. Observe that

$$
\frac{(x-3)^2}{2(x+3)} = \frac{x-9}{2} + \frac{36}{2x+6}.
$$

Since $x$ is odd the first term is an integer so we just require
$2x+6$ to be a factor of 36. Thus $2x+6 \in
\{ 1,2,3,4,6,9,12,18,36 \}$. Since $2x+6$ is even we can eliminate any odd
candidates, and since $x\geq0$ we have $2x+6\geq 6$. Thus
$2x+6 \in \{ 6,12,18,36\}$. So we solve to see that
$x=0,3,6,15$. $\blacksquare$
\end{solution}
\vspace{0.5cm}
\newpage

\begin{problem}{Problem 6}
Prove the Power of a point theorem.
\end{problem}

\begin{solution}{Solution to problem 6}
\textbf{(a):} Consider constructing $\overline{DA}$ and $\overline{BC}$,

\begin{center}
\begin{asy}
size(120);
pair A = (-4,3);
pair C = (4,3);
pair B = (0,5);
pair D = (0,-5);
pair P = intersectionpoint(A--C, B--D);
pen pastel_blue = rgb(0.68, 0.85, 0.90);
pen pastel_pink = rgb(0.99, 0.71, 0.75);
pen line_pen = black + linewidth(0.75);
fill(circle((0,0),5), pastel_blue + opacity(0.2));
fill(A--P--D--cycle, pastel_blue + opacity(0.2));
fill(B--P--C--cycle, pastel_pink + opacity(0.2));
draw(circle((0,0),5), blue);
draw(A--C, red);
draw(B--D, red);
draw(A--D, dashed + red);
draw(B--C, dashed + red);
label("$A$", A, W);
label("$C$", C, E);
label("$B$", B, N);
label("$D$", D, S);
label("$P$", P, E);
\end{asy}
\end{center}


Due to the Inscribed angle theorem $\angle ADP = \angle BCP$, and since
opposite angles are equal $\angle APD = \angle BPC$. Since there are two
equal angles and angles in a triangle add to $180^{\circ}$ we also have
$\angle DAP= \angle PBC$. Now since there are three equal angles the
triangles $\triangle DAP$ and $\triangle PBC$ are similar. \vspace{0.2cm}

Thus $AP/DP = BP/CP$ which
re-arranges to $AP\cdot CP = BP \cdot DP$ .
\\

\textbf{(b):} Consider constructing the triangles $\triangle ABP$ and
$\triangle DCP$.

\begin{center}
\begin{asy}
import graph;
size(120);

real labelscalefactor = 1.0;

// Image dimensions
real xmin = -13.91341807885702, xmax = 15.180591129029986, ymin = -7.699509283857027, ymax = 7.771251695529754;

/* --- Define Points (Pairs) --- */
// These coordinates are taken from your draw/label commands
pair A = (-3.489002731044119,-1.9562361674313966);
pair D = (0.5974092528636418,3.9551361777558185);
pair P = (7.,-4.);
pair B = (3.9914183964301753,-0.26187627735013763);
pair C = (2.499312903336305,-3.123048992765987);

/* --- Define Colors & Pens --- */
pen pastel_blue = rgb(0.68, 0.85, 0.90);
pen pastel_pink = rgb(0.99, 0.71, 0.75); // <-- New color for polygon fill

/* --- Draw Figures with Colors --- */
// Circle
path circ = circle((0.,0.), 4.);
fill(circ, pastel_blue + opacity(0.2)); // Fill the circle

// *** FILL POLYGONS HERE ***
// Define the polygons as paths
path poly_ABP = A--B--P--cycle;
path poly_DCP = D--C--P--cycle;

// Fill the polygons (with your chosen color and opacity)
fill(poly_ABP, blue + opacity(0.1));
fill(poly_DCP, pastel_pink + opacity(0.2));

/* --- Draw Lines and Circle Border --- */
// (Drawn *after* filling so they are on top)
draw(circ, blue); // Circle border

// Solid lines
draw(A--P, red);
draw(D--P, red);

// Dashed lines
draw(A--B, dashed + red);
draw(D--C, dashed + red);

/* --- Dots and Labels --- */
// Using the pair variables makes this cleaner
label("$A$", A, NE * labelscalefactor);
label("$D$", D, NE * labelscalefactor);
label("$P$", P, NE * labelscalefactor);
label("$C$", C, NE * labelscalefactor);
label("$B$", B, NE * labelscalefactor);
clip((xmin,ymin)--(xmin,ymax)--(xmax,ymax)--(xmax,ymin)--cycle);
\end{asy}
\end{center}

Again due to the Inscribed angle theorem $\angle PAB = \angle CDP$.
Since $BP$ and $CP$ are common to both triangles
we deduce that $\angle APB = \angle DPC$. Since there are two equal angles
and angles in a triangle add up to $180^\circ$ we also have
$\angle ABP = \angle PCD$. Now since there are three equal angles the
triangles $\triangle ABP$ and $\triangle DCP$ are similar. \vspace{0.2cm}

Thus $AP/BP = DP/CP$ which
re-arranges to $AP\cdot CP = BP \cdot
DP$.\\

\textbf{(c):} Consider constructing the triangles $\triangle ABP$ and
$\triangle BCP$.

\begin{center}
\begin{asy}
/* Geogebra to Asymptote conversion, documentation at artofproblemsolving.com/Wiki go to User:Azjps/geogebra */
import graph; size(120);
real labelscalefactor = 1.0; /* changes label-to-point distance */
real xmin = -9.469163778645866, xmax = 9.583363840658867, ymin = -7.075242178993681, ymax = 5.837817387494271;  /* image dimensions */
pen ttttff = rgb(0.2,0.2,1.); pen ccqqqq = rgb(0.8,0.,0.); pen yqqqqq = rgb(0.5019607843137255,0.,0.); pen zzttqq = rgb(0.6,0.2,0.); pen ttffcc = rgb(0.2,1.,0.8);
pair B = (0.,-5.), P = (8.,-5.), A = (3.3184235880875335,3.74006215055636), C = (4.95215586502357,0.6900378891862754);
pen pastel_blue = rgb(0.68, 0.85, 0.90);

filldraw(circle((0.,0.), 5.), pastel_blue + opacity(0.2), ttttff);
filldraw(arc(B,0.8958126070288543,359.42543793229623,408.44701114793276)--(0.,-5.)--cycle, yqqqqq + opacity(0.10000000149011612), linetype("4 4") + yqqqqq);
filldraw(arc(A,0.8958126070288543,-110.7907497333701,-62.3437385854374)--(3.3184235880875335,3.74006215055636)--cycle, yqqqqq + opacity(0.10000000149011612), linetype("4 4") + ccqqqq);

filldraw(B--P--C--cycle, zzttqq + opacity(0.10000000149011612), zzttqq +
opacity(0.10000000149011612));
filldraw(A--B--P--cycle, ttffcc + opacity(0.10000000149011612), ttffcc +
opacity(0.10000000149011612));
/* draw figures */
draw(B--P, ccqqqq);
draw(A--P, ccqqqq);
draw(B--C, dashed + ccqqqq);
draw(A--B, dashed + ccqqqq);
/* dots and labels */
label("$B$", (0.08723215613926336,-4.847961214935062), NE * labelscalefactor);
label("$P$", (8.070211666712167,-4.808924640311234), NE * labelscalefactor);
label("$A$", (3.4053409991646757,3.9352680754263316), NE * labelscalefactor);
label("$C$", (5.025358846053554,0.8513786801438865), NE * labelscalefactor);
clip((xmin,ymin)--(xmin,ymax)--(xmax,ymax)--(xmax,ymin)--cycle);
/* end of picture */
\end{asy}
\end{center}


Now clearly $\angle BPC = \angle APB$ because $BP$ and
${CP}$ are common to both triangles. Then also
$\angle PBC = \angle BAP$ due to the Alternate segment theorem.
Since there are two equal angles
and angles in a triangle add up to $180^\circ$ we also have
$\angle ABP = \angle BCP$. Since we have three equal angles
triangles $\triangle ABP$ and $\triangle BCP$ are similar.
\vspace{0.2cm}

Thus ${AP}/{BP} = {BP}/{CP}$ which
re-arranges to ${AP}\cdot {CP} = {BP}^2$.
$\blacksquare$
\end{solution}
\vspace{0.5cm}
\newpage

\begin{problem}{Problem 7}
Consider the thirteen factorials $0!, 1!, 2!, \ldots , 12!$. We want to divide these up into three sets of size 4
and one set of size 1. Moreover, we would like the product of the elements in each of the sets of size 4
to be a square number.
\vspace{0.2cm}

In how many ways can this be achieved?
\end{problem}

\begin{solution}{Solution to problem 7}
The following table will be of use to us, even numbers are in red whilst
odd numbers are in blue.


\[
\begin{tblr}{
  colspec={c|*{13}{c}},
  hlines,
  vlines,
}
n & 0 & 1 & 2 & 3 & 4 & 5 & 6 & 7 & 8 & 9 & 10 & 11 & 12 \\ \hline
v_2(n!)  & \textcolor{red}{0} & \textcolor{red}{0} & \textcolor{blue}{1} & \textcolor{blue}{1} & \textcolor{blue}{3} & \textcolor{blue}{3} & \textcolor{red}{4} & \textcolor{red}{4} & \textcolor{blue}{7} & \textcolor{blue}{7} & \textcolor{red}{8} & \textcolor{red}{8} & \textcolor{red}{10} \\
v_3(n!)  & \textcolor{red}{0} & \textcolor{red}{0} & \textcolor{red}{0} & \textcolor{blue}{1} & \textcolor{blue}{1} & \textcolor{blue}{1} & \textcolor{red}{2} & \textcolor{red}{2} & \textcolor{red}{2} & \textcolor{red}{4} & \textcolor{red}{4} & \textcolor{red}{4} & \textcolor{blue}{5} \\
v_5(n!)  & \textcolor{red}{0} & \textcolor{red}{0} & \textcolor{red}{0} & \textcolor{red}{0} & \textcolor{red}{0} & \textcolor{blue}{1} & \textcolor{blue}{1} & \textcolor{blue}{1} & \textcolor{blue}{1} & \textcolor{blue}{1} & \textcolor{red}{2} & \textcolor{red}{2} & \textcolor{red}{2} \\
v_7(n!)  & \textcolor{red}{0} & \textcolor{red}{0} & \textcolor{red}{0} & \textcolor{red}{0} & \textcolor{red}{0} & \textcolor{red}{0} & \textcolor{red}{0} & \textcolor{blue}{1} & \textcolor{blue}{1} & \textcolor{blue}{1} & \textcolor{blue}{1} & \textcolor{blue}{1} & \textcolor{blue}{1} \\
v_{11}(n!) & \textcolor{red}{0} & \textcolor{red}{0} & \textcolor{red}{0} & \textcolor{red}{0} & \textcolor{red}{0} & \textcolor{red}{0} & \textcolor{red}{0} & \textcolor{red}{0} & \textcolor{red}{0} & \textcolor{red}{0} & \textcolor{red}{0} & \textcolor{blue}{1} & \textcolor{blue}{1} \\
\end{tblr}
\]

We define the \textit{score} of a set, to be the product of the elements in
that set. \vspace{.4cm}

\textit{Claim 1: The number in the set of size 1 must be $6!$.}
\vspace{.2cm}

\textit{Proof:}
Suppose the three sets of size 4 have scores $A^2, B^2$ and $C^2$. Then
suppose the number we place in the set of size 1 is $k!$. We observe that
$A^2 B^2 C^2 k! = 0! \cdot 1! \cdot 2! \cdots 12!$.
Let $N = (0! \cdot 1! \cdot 2! \cdots 12!)/k! = (ABC)^2$. Since
$v_5(0! \cdot 1! \cdot 2! \cdots 12!) = 11$, $k!$ must have an odd 5-adic
valuation which means $k \in \{ 5,6,7,8,9 \}$, \vspace{0.2cm}

If $k = 7,8,9$ then $v_7(N) = 5$ so $N$ isn't square.
If $k = 6$ then $N$ is in fact a square.
If $k = 5$ then $v_3(N) = 25$ so $N$ isn't square.
Thus $k = 6$. $\square$ \vspace{0.4cm}

\textit{Claim 2: Let $p$ be some prime. There must be an even number of
factorials that have an odd $p$-adic valuation
in each set of size 4, similarly there must be an even number of factorials
that have an even $p$-adic valuation in each set of size 4.}
\vspace{0.2cm}

\textit{Proof:} Suppose there are $x$ factorials that have an even $p$-adic
valuation and $y$ factorials that have an odd
$p$-adic valuation in one set of size 4. Let $S$ be the score of that set.
Then $v_p (S) \equiv y \equiv 0 \ \text{(mod } 2)$ since
perfect squares only have even $p$-adic valuations. The claim follows from
the fact that y is even. $\square$
\vspace{0.4cm}

\textit{Claim 3: $11!$ and $12!$ must be in the same set. }
\vspace{0.2cm}

\textit{Proof:} Since $11!$ and $12!$ are the only numbers with factors of
11, and they both have exactly one factor of 11, they must be in the same
set otherwise one of the sets would have a score with an
odd 11-adic valuation, but since the scores are all perfect squares this
is a contradiction. $\square$
\vspace{0.4cm}

\textit{Claim 4: The set that contains $11!$ and $12!$ is either
$\{2!, 3!, 11!, 12! \}$ or $\{2!, 4!, 11!, 12!\}$.}
\vspace{0.2cm}

\textit{Proof:}
Now in the set that contains $11!, 12!$ the two other numbers $x!$ and
$y!$ must satisfy $v_3(x!\cdot y!) \equiv 1 \ \text{(mod }2)$ because
$v_3(11!\cdot 12!) = 9$ is odd and $x!\cdot y!\cdot 11!\cdot 12!$ is a
perfect square. Call a number 3-\textit{blue} if it has an odd 3-adic
valuation and is not one of $11!, 12!$. Call a number 3-\textit{red} if it
has an even 3-adic valuation and is not one of $11!, 12!$.\vspace{0.2cm}

Out of $x!$ and $y!$ exactly one of them must be 3-blue because otherwise
$x!\cdot y!$ would have an even 3-adic valuation. The only 3-blue numbers
are $3!,4!$ and $5!$ and they all have an odd 2-adic valuation. Suppose
WLOG that $x!$ is 3-blue. Then $y!$ must have an odd 2-adic valuation
also because $v_2(11!\cdot 12!\cdot x!)$ is odd and
$x!\cdot y!\cdot 11!\cdot 12!$ is a perfect square. We can use our table
and the information deduced thus far to see that $y! \in \{ 2!, 8!, 9! \}$.
However $y! \neq 8!, 9!$ because then $v_5(x!\cdot y!\cdot 11!\cdot 12!)=3$
which is odd, but this is a contradiction. Thus $y!=2!$. Now it can be
checked that the set that contains $11!, 12!$ must be either
$\{2!, 3!, 11!, 12 \}$ or $\{2!, 4!, 11!, 12!\}$. $\square$
\vspace{0.4cm}

Now suppose that we have placed $6!$ in the set of size 1, and we have
selected the numbers that will be in the set that contains $11!$ and
$12!$ using claim 4. The only remaining numbers that have factors of 7
are $7!, 8!, 9!$ and $10!$. Let's call these numbers \textit{elegant}.
Notice that elegant numbers all have exactly one factor of 7.
Now due to claim 2 each set must contain an even number of elegant numbers.
It is important to note that we can't have all four elegant numbers in one
set since $v_3(7!\cdot 8!\cdot 9!\cdot 10!) = 5$ is odd and so the product
of the four elegant numbers isn't a perfect square.
\vspace{0.2cm}

One of the sets will look like this $\{ 10!, x!, \_, \_ \}$ where $x!$ is
an elegant number other than $10!$. Since $v_5(10!\cdot x!) = 3$, $5!$ must
be in this set because it is the only remaining number with a factor of 5.
Thus that set looks like $\{ 10!, x!, 5!, \_ \}$. However now using the
table it can be checked that $v_3(10!\cdot x!\cdot 5!)$ is odd and so the
final number represented by \_ must have an odd 3-adic valuation. That
means it must be $3!$ or $4!$, whichever one wasn't used in the set
that contains $11!$ and $12!$. \vspace{0.2cm}

That means the third and final set must look like $\{ 0!, 1!, y!, z!\}$
where $y!$ and $z!$ are the two remaining unused elegant numbers.
One can use the table to check that only $(y!, z!) = (8!, 9!)$ works.
This means that $x! = 7!$.
\vspace{0.2cm}


Now putting all of this together our sets look like so, let $a!, b!$ be
$3!, 4!$ in some order,
$\{6!\}$, $\{0!,1!,8!,9!\}$, $\{2!, a!, 11!, 12!\}$ and
$\{b!, 5!, 7!, 10!\}$. Now it can be checked that, no matter the
order of $a!, b!$, each set has a score that is a perfect square. There
are 2 ways to choose $(a!, b!)$ so the final answer is 2. $\blacksquare$
\\

\textit{Comments: This problem is arguably the hardest on the sheet,
the solution I found involves some facts from number theory and a lot of
patience to work through all of the cases.}
\end{solution}
\vspace{0.5cm}

\newpage
\begin{problem}{Problem 8}
Find all sets of three consecutive positive integers such that each is either a prime or a power of a prime.
(We do not regard 1 as a power of a prime in this question.)
\end{problem}

\begin{solution}{Solution to problem 8}
Call a set $\textit{good}$ if it contains three consecutive positive
integers such that they are all prime or powers of a prime.
Let $n \in \mathbb{N}$ such that $n, n+1$ and $n+2$ are all primes or powers
of a prime. \vspace{0.2cm}

If $n$ is even then $n=2$ since 2 is the only even prime.
Also since $n+1 = 3$ and $n+2 = 2^2$ the set $\{2, 3, 4\}$ is a good set.
\vspace{0.2cm}

If $n$ is odd then $n+1$ is even which means it must be of the form $2^k$
for $k \in \mathbb{N}$. Similarly $n = 2^k-1$ must be of the form $p^a$
and $n+2=2^k+1$ must be of the form $q^b$ for odd primes $p,q$ and
$a,b \in \mathbb{N}$. Now we notice that
$n = 2^k - 1 \equiv (-1)^k - 1 \ \text{(mod } 3)$ and
$n + 2 = 2^k + 1 \equiv (-1)^k + 1 \ \text{(mod } 3)$. Thus
if $k$ is even then $3 \mid n$, and if $k$ is odd then $3 \mid (n+2)$.
\vspace{0.2cm}

Suppose that $k$ is even, then we have $n = 3^a$ and $n+1 = 2^k$. We deduce that
$2^k - 3^a = 1$ or equivalently $(2^{k/2}-1)(2^{k/2}+1)=3^a = n$. Now we
have two
powers of 3 that are 2 apart, this only happens for $(3^0, 3^1)$ because
powers of 3 grow too quickly past this point. So we have $2^{k/2}-1=1$ and
$2^{k/2}+1=3$ which gives $a=1, k=2$ and so $n=3$. Also since $n+1=2^2$ and
$n+2=5$ the set will be good. So the only good set in this
case is $\{3, 4, 5 \}$.
\vspace{0.2cm}

Suppose that $k$ is odd, then $n+2=3^b$ and $n+1 = 2^k$. This time we have
$3^b - 2^k = 1$. We use LTE to finish. We need that $v_2(3^b-1^b) = k$,
there are two cases. If $b$ is odd then $v_2(3^b-1^b) = v_2(3-1) = 1 = k$.
We find that $k=1, b=1$ but that implies that $n = 1$ which is not a power
of a prime in the context of this question. \vspace{0.2cm}

Thus $b$ must be even and we
have $v_2(3^b-1^b)= v_2(3-1) + v_2(3+1) + v_2(b) - 1 = k$ which re-writes
as $v_2(b) = k-2$. So we find that $2^{k-2} \mid b$ and so $b \geq 2^{k-2}$.
Clearly this inequality will be very restrictive since

$$
1 = 3^b - 2^k \geq 3^{2^{k-2}} - 2^k
$$

will only hold for $k \leq 3$ since that RHS grows too quickly. We have
already eliminated the case $k=1$ and it can be checked that $k=3, b=2$
is a valid solution here. Thus $n = 7$, and since $n+1=2^3$ and $n+2 = 3^2$
the set $\{7,8,9 \}$ is the only good set here. \vspace{0.2cm}

Putting all of that together the only good sets are $\{2,3,4 \}$,
$\{ 3,4,5 \}$ and $\{7,8,9 \}$. $\blacksquare$
\\

\textit{Comments: This one was fun! The case when $k$ is odd is probably
the most challenging part of the problem, I personally can never remember
all of the edge cases of LTE so I had to look it up. Of course my arguements
that say something like "...grows too quickly" are not at all rigorous but
I omit
a formal justification because it'd make the solution really long, but I
don't mind adding justification if you want!}
\end{solution}

\end{document}
